% 第四问（基于 paper_q3_improvements 的隔离优化方案）
\section{问题四：波动电价下的滚动储能调度}

\subsection{建模思路}
问题四在问题三的滚动承诺框架上引入附件4的时变电价，并采用队友提出的负荷偏差修正方法。对于每个发布时刻，首先使用截至该时刻可获得的历史数据预测未来负荷和光伏；随后从最近30个有效历史日提取同一发布时刻的整段负荷--光伏残差，构造情景集并求解线性规划。预测阶段不读取目标时段的真实数据，真实数据仅在计划冻结后的回放结算中使用。

在日内更新时刻 $k\in\{36,72,108\}$（分别对应 $6{:}00,12{:}00,18{:}00$）上，负荷预测采用最近12个十分钟残差的中位数作为偏差估计，并令其随预测步长衰减；$0{:}00$ 的日初计划使用未修正的基础预测：
\[
\widehat L_{d,k+h}=\max\left\{0,\widehat L^{(0)}_{d,k+h}+b_{d,k}\exp(-h/36)\right\},
\quad
b_{d,k}=\operatorname{median}_{j=k-12}^{k-1}(L_{d,j}-\widehat L^{(0)}_{d,j}).
\]
光伏预测沿用基础预测器，误差场景由最近30日同一发布时刻的整段残差生成。计划电价采用目标日期前7日同一时段价格，真实波动电价只用于事后结算，从而避免未来价格信息泄漏。

\subsection{两种执行规则}
问题4-2在每日 $0{:}00$ 形成全天计划并锁定；问题4-3在四个发布时刻滚动更新，但每次仅执行接下来6小时。更新时继承已经生效的承诺量，对尚未执行区间重新优化。上调量和取消量分别按每次更新的绝对变化累计结算，不允许用最终净调整抵消中间交易。实际运行中，以真实负荷、真实光伏和储能状态执行已锁定承诺，若仍有缺口则按 $5p_t$ 购买紧急电量。日末实际储能量连续传递至下一天。

\subsection{全年结果}
采用 P1（2小时负荷偏差修正）方案，覆盖2025年2月1日至12月31日共334天、48,096个十分钟时段。结果如表~\ref{tab:q4p1}所示。问题4-2全年总费用为 $14,102,925.03$ 元，其中紧急购电费为 $612,543.66$ 元；引入每6小时滚动更新后，问题4-3总费用降至 $13,790,040.90$ 元，紧急购电费降至 $81,492.29$ 元。日内更新使总费用下降约 $2.22\%$，紧急购电费用下降约 $86.69\%$，说明最新负荷信息主要改善了短时供能缺口，而非仅改变跨日套利。

\begin{table}[H]
\centering
\caption{基于队友Q3改进方案的第四问全年结果}
\label{tab:q4p1}
\begin{tabular}{lrr}
\toprule
指标 & 问题4-2（每日决策） & 问题4-3（6小时滚动）\\
\midrule
总费用/元 & 14102925.03 & 13790040.90\\
紧急购电费用/元 & 612543.66 & 81492.29\\
运行天数 & 334 & 334\\
逐时记录数 & 48096 & 48096\\
\bottomrule
\end{tabular}
\end{table}

\subsection{数值检验与信息边界}
逐时结果满足功率平衡和储能状态递推，最大平衡残差与状态残差均约为 $10^{-12}$ kWh；所有实际储能量均位于 $[1200,10800]$ kWh。结果文件同时给出计划购电、最终承诺、上调、取消、紧急购电和弃电字段，可逐项复算每日结算费用。需要强调的是，上述结果是“预测--锁定--真实回放”的因果模拟，不是把全年真实负荷、光伏和电价提前交给优化器得到的完全信息下界。

\subsection{结论}
在同一基础预测和场景生成框架下，问题4-3通过日内更新及最近2小时负荷偏差修正，取得比每日一次决策更低的回放费用，尤其显著降低紧急购电风险。两种改变共同作用，因此本实验不能单独量化偏差修正的贡献。若题目强调实际运行可靠性，推荐采用问题4-3；问题4-2可作为计算资源有限时的简化备选。
